\documentclass[11pt,a4paper]{article}
\usepackage[utf8]{inputenc}
\usepackage[margin=1in]{geometry}
\usepackage{amsmath,amssymb}
\usepackage{xcolor}
\usepackage{xr-hyper}
\externaldocument{main} % reads main.aux; omit the .tex extension
\usepackage{hyperref}   % only if not already loaded
% \usepackage{hyperref}

% Custom styling for comments and changes
\definecolor{refcolor}{RGB}{0, 51, 102}
\definecolor{quotecolor}{RGB}{80, 80, 80}

\newcommand{\refcomment}[1]{%
    \medskip\noindent\textbf{\color{refcolor}Comment:}\ %
    {\itshape #1}\par\smallskip
}

\newcommand{\response}{\noindent\textbf{Response:}\ }

\newcommand{\addedtext}[1]{%
    \smallskip\begin{quote}\color{quotecolor}\small #1\end{quote}\smallskip
}

\title{\textbf{Response to Referee Reports}\\
\large JHEP Manuscript ID: JHEP\_027P\_0726}
\author{\textbf{Authors:} Austin Townsend, Marc Osherson, Mike Hildreth, Stefano Castruccio}
\date{\today}

\begin{document}

\maketitle

\noindent Dear Editor and Referees,

\medskip
We thank both referees for their careful reading and constructive suggestions. Their comments have helped us clarify the statistical motivation for exponential mixtures and better distinguish the conclusions supported by our studies from questions requiring further investigation.

Below, we include a summary of changes and then respond to each comment. Where we have not undertaken a suggested additional study, we explain our reasoning.

\section*{Summary of Changes}

We have substantially revised the presentation of the model’s statistical motivation, clarified the fitting and simulation procedures, and expanded the discussion of limitations. The principal changes are summarized below, with detailed explanations provided in the point-by-point responses.

1. Abstract and introduction: We now define the finite exponential mixture explicitly as a weighted sum of exponential densities with non-negative weights and identify the quoted dataset sizes as observation counts. The introduction acknowledges previous use of exponential sums in background modeling and distinguishes the statistical motivation for the proposed construction from its empirical validation.

2. GPD motivation and mixture construction (Section 2): We have rewritten the development of the model to proceed directly from the positive-shape generalized Pareto distribution (GPD), through its exact representation as an exponential mixture with gamma-distributed rates, to a general mixing distribution and its finite approximation. We have added the exponential case to the GPD definition and clarified the allowed mass ranges, log-convexity and log-concavity, and the distinction between the limiting behavior described by extreme value theory and the behavior over a finite analysis interval.

3. Complete monotonicity and approximation properties (Section 2): We have made the approximation statement precise by specifying restriction and normalization over a common finite analysis interval. We have replaced the discussion of the Normal-distribution tail with examples of completely monotone functions, including power laws, shifted power laws, and stretched exponentials, and explained how closure under addition and multiplication broadens the range of shapes represented by the mixture family.

4. Optimization and model selection (Section 3): We now specify the accepted minimizer statuses, the selection of the lowest negative log-likelihood among accepted fits, and the distinction between fits to independent pseudo-datasets and fits to bootstrap resamples. The dataset-specific restart and fit-call budgets are stated explicitly and identified as heuristic choices that do not certify global optimality. We also express AIC and BIC as likelihood-improvement thresholds and explain why they are used as practical selection rules in this setting, where standard regularity assumptions may fail.

5. Likelihood definitions and fit comparisons (Section 3; Tables 1–2): The fit-comparison tables have been updated to include the degrees of freedom and selected mixture orders. Figures 2 and 4 now use linear mass axes with logarithmic vertical axes, and Figures 1 and 3 have been replaced with vector versions. We have revised captions and equation references, grouped the four diphoton reference functions under Eq. (3.4), and corrected typographical and terminology inconsistencies. 

6. Bias and coverage (Section 4): The coverage discussion now explains weak identifiability through nearly coincident rates and negligible component weights. We describe the bootstrap resampling, initialization, repeated model selection, and percentile-interval construction explicitly, and distinguish these steps from the evaluation of coverage across independent pseudo-datasets. We also quantify the statistical precision of the coverage estimates and acknowledge correlations across mass values and undercoverage in the tail.

7. Spurious-signal assessment (Section 4; Figure 7): We have replaced the Hessian-standardized signal yield with the signed profile-likelihood statistic \(r_0=\operatorname{sign}(\hat{\mu})\sqrt{-2\ln\lambda(0)}\), explaining its relation to the usual standardized estimator under the Wald approximation. The signal-plus-background model is now written in terms of a signal fraction, with fixed signal mean and width and negative fractions permitted subject to non-negativity of the total density. We specify that the background-only selected mixture order is held fixed during profiling, and have updated Figure 7 and its caption to identify the statistic and the central intervals across pseudo-datasets.

8. Discussion and conclusions (Sections 5–6): We have substantially reorganized the discussion to distinguish the evidence supporting density estimation from the additional validation required for resonance searches. The revised text explicitly addresses possible signal absorption, the need to assess power alongside false-positive rates, and the limited sample-size range of the simulation studies. Signal-injection, sensitivity, and additional sample-size studies remain outside the present revision and are identified as outstanding work. We have also expanded the discussion of optimization, adaptive stopping criteria, ordered-rate parameterizations, regularization, and endpoint modeling, and qualified statements about protection against overfitting.

\bigskip
\hrule
\bigskip

\section*{Response to Referee 1}

\refcomment{Abstract "motivate the finite exponential mixture as a" This phrase is not
self-explanatory. I was imagining some sort of e*a + bx + cx2. Googling finds your article as top hit, so it does not seem to be in common use.}
\response We agree that the abstract can be clearer. This sentence has been amended:
\addedtext{In this work, we investigate the finite exponential mixture, defined as a weighted sum of exponential densities with non-negative weights, as a flexible semi-parametric class of functions for approximating falling distributions.}

\refcomment{I misinterpreted ’n= 28, 619, 185 and n = 5, 036’ as somehow identifying
numbers of a dataset, not the size.}
\response Modified text:
\addedtext{Using two published datasets containing 28,619,185 and 5,036 observations, respectively, ...}

\refcomment{Section 2 starts with pure motivations, but quickly reverts into an ad-hoc
model. But having explained the benefits of the GPD in including models with upper cut-offs, these are not supported.}
\response We agree that the distinction between endpoint behavior and the shape over a finite analysis range requires clarification. The proposed mixture does not describe the vanishing spectrum at a finite kinematic endpoint. Instead, we use the positive-shape GPD as a reference family whose power-law behavior is compatible with spectra over the intervals considered here, and generalize its distribution of exponential rates. We have revised Section 2 to clarify this motivation and discuss the separate treatment of endpoint effects in Section 5. See added text:

\addedtext{In collider applications, the physical upper bound at the center-of-mass energy motivates a negative-shape GPD as a possible description sufficiently close to the endpoint. The range over which endpoint effects become important depends on the spectrum: they may be confined to a narrow interval near the kinematic limit or extend to substantially lower masses. Close to the endpoint, the vanishing spectrum also makes observations increasingly rare. The practical importance of modeling these effects therefore depends on both their extent across the fitted range and the available event count. Over the mass intervals considered by the analyses presented here, the spectra are instead predominantly log-convex, favoring the positive-shape GPD as an alternative reference shape whose power-law behavior is compatible with established parameterizations of non-resonant backgrounds. This compatibility motivates the positive-shape GPD as a starting point, although its specific functional form may be too restrictive to capture variations across the full analysis range. Its representation as an exponential mixture with gamma-distributed rates then suggests a systematic generalization: relaxing the gamma restriction allows for more flexibility.}

\refcomment{Eqn 2.4 is rather straightforward and its not clear what is gained by passing
through the continuous eqn 2.2}
\response The purpose is to show that in going from the continuous to the finite case we are making a particular approximation. We can make this more clear by moving from the exponential-gamma representation of the positive shape GPD to the finite exponential mixture. This section has been amended to include the following text:
\addedtext{For the following mixture representations, we measure $x$ relative to the lower threshold, so that $x$ denotes the excess above $u$ and the threshold is located at zero. For $\xi>0$, the GPD has the exponential-mixture representation
$$
f_{\text{GPD}}(x\mid\sigma,\xi>0)
=\int_0^\infty \beta e^{-\beta x}
\pi_\Gamma(\beta\mid a,b) d\beta,
$$
where
$$
\pi_\Gamma(\beta\mid a,b)
=\frac{b^a}{\Gamma(a)}\beta^{a-1}e^{-b\beta},
\qquad
a=\frac{1}{\xi},\qquad b=\frac{\sigma}{\xi},
$$
where $a$ and $b$ are the shape and rate parameters of the Gamma distribution. The exponential case $\xi=0$ corresponds to a rate distribution concentrated at $\beta=1/\sigma$.

This identity motivates replacing the gamma distribution by a general probability distribution $G$ over positive rates:
$$
f_G(x)=\int_0^\infty \beta e^{-\beta x} dG(\beta).
$$
The resulting family contains the positive-shape GPD and exponential reference models while allowing more flexible shapes. Approximating $G$ by a discrete distribution gives the finite exponential mixture density:
$$
f_k(x)=\sum_{i=1}^{k}w_i\beta_i e^{-\beta_i x},
\qquad
w_i\geq0,\quad \sum_{i=1}^{k}w_i=1,\quad \beta_i>0.
$$
The set of rates $\beta_i$ and the corresponding weights $w_i$ provide a finite approximation of the mixing distribution $G$ and likely lack any physical interpretation alone.}

\refcomment{"In some regimes the Normal distribution tail converges to the GPD" I am
not clear what we learn from this - you moved away from the GPD, right?}
\response We have not moved away from the GPD, the exponential mixture is a generalization of the unbounded GPD. This statement is used to give the reader an idea of the generality of EVT. This is somewhat vague though and potentially misleading, the Normal distribution converges to an exponential and the rate of this convergence is famously slow. This and the preceding sentences discussing the Normal distribution have been removed and the paragraph has been amended to contain examples of completely monotone functions instead of counter-examples. Here is the amended paragraph:

\addedtext{Any completely monotone function on $(0,\infty)$ admits a representation as a positive mixture of exponential functions \cite{bernstein} and can be uniformly approximated by finite positive exponential sums on compact intervals contained in $(0,\infty)$ \cite{McGlinn_1978}. Thus, after restricting and normalizing both the target function and the mixture over the same finite analysis interval, finite exponential mixtures can approximate arbitrarily closely any density obtained from a completely monotone function. Familiar examples of completely monotone functions include decreasing exponentials $e^{-\beta x}$, power laws $x^{-\alpha}$, shifted power laws $(1+x/\sigma)^{-\alpha}$, and stretched exponentials $e^{-(x/\sigma)^q}$, for $\alpha,\beta,\sigma>0$ and $0<q\leq1$, on $x>0$. The class of completely monotone functions is closed under addition and multiplication, so positive sums of these functions and products such as exponentially damped power laws also belong to the class. Exponential mixtures therefore provide a common representation for several familiar falling shapes without requiring a particular combination of factors to be specified in advance.}

\refcomment{At the top of p5 you make it clear that the convergence is not straight-
forward ’multiple random restarts’. This sounds slow and cumbersome. You need a procedure to define success...what is it?}
\response We agree that the details here are insufficient. We have modified the text to be more clear:
\addedtext{Numerical optimization is performed with \texttt{MINUIT} and the \texttt{MIGRAD} algorithm~\cite{Minuit2}. Because the likelihood is non-convex in the mixture parameters, each model is fitted from multiple randomized starting points. At each starting point, minimization is repeated until the returned status is 0 (successful minimization), 1 (the covariance matrix was forced positive definite), or 2 (the Hessian calculation was invalid), or until a prescribed limit on the number of fit calls is reached. Among the accepted results, the parameter values yielding the lowest negative log-likelihood are retained.
}

\refcomment{You introduce AIC/BIC. It is very common in the field to use Wilkes tests, or the derivative 2*delta(ll) < X when testing different numbers of parameters. Did you have a good reason to use a different approach?}
\response AIC and BIC offer practical ways to balance improvement in the fit against additional parameters. They are similar to the likelihood-improvement threshold suggested by the referee: AIC accepts the larger model when $2\Delta\ell>2\Delta p$, while BIC requires $2\Delta\ell>\Delta p\log n$, where $\Delta p$ is the number of additional parameters and $n$ is the event count. Adding one exponential component introduces two free parameters, giving thresholds of $4$ and $2\log n$, respectively.

For finite mixtures, violations of the regularity conditions cast doubt on the usual asymptotic formulae used to interpret likelihood-ratio tests. Such tests may nevertheless provide useful selection rules, but their performance in this setting requires further study. Bootstrap calibration is an alternative, although repeating it for every model-selection decision is computationally prohibitive for our studies. We also acknowledge that the assumptions underlying the usual theoretical justification of AIC and BIC may be violated. We therefore use these criteria as practical choices and evaluate their performance through the simulation studies, rather than relying solely on their asymptotic justification.


\refcomment{on p6 you write 80 restarts and 100 retries ’since the dataset is large’. This does not seem like an automated procedure.}
\response We agree that the choices of restart and retry budgets are heuristic and that an automated stopping criteria would be a nice improvement. The revised discussion identifies the development and validation of adaptive stopping criteria and more efficient optimization strategies as future work. Here is the modified discussion:
\addedtext{Optimization remains a practical limitation of the present implementation. The likelihood is non-convex in the mixture parameters and can contain nearly flat directions when components have similar rates or negligible weights. Random restarts reduce dependence on initialization, but the finite search used here does not certify that the global maximum has been reached. The restart and retry budgets are heuristic, and their sufficiency across datasets and model complexities has not been established. Residual optimization error can also affect model selection when competing information-criterion values are close. An adaptive stopping criterion could help allocate computational effort according to the observed progress of the optimization, reducing reliance on manually chosen budgets while providing a reproducible rule for terminating the search. Beyond refining the current restart procedure, alternative strategies such as split-and-merge methods developed for Gaussian mixtures \cite{split-and-merge} warrant investigation for their potential to improve optimization reliability and efficiency in exponential mixtures.}

\refcomment{table 1/2: why are you not reporting the Ndf of the fit and order of
the Exponential mixture? Its not clear why you quote the chi2 here at
all - the ll does not need this arbitrary rebinning to have 30 events and is
just as informative for comparing values. Indeed for the diphoton dataset the visually obvious deviations are above 2000 and are perhaps lost in the rebinning.}
\response The chi2 is commonly used in HEP and we include it as a familiar diagnostic. We agree that merging low-count bins reduces resolution in the tail and that the resulting statistic should not replace inspection of the residuals. We have added the degrees of freedom and selected mixture order to the tables and captions.

\refcomment{page 10 "Since fits to exponential mixtures of varying complexity may
have weakly identified parameters," What do you mean? Anticorrelated
parameters may give unstable results? This is a worrying flag that it would be
good to explain a bit more. The impact is that you try to estimate coverage
of the simulation using bootstrap on the data. This is a fine approach but
does but measures not true density but fluctuations of the observed data,
which is not quite the same thing - if your dataset is ’up’ in one place the
sampling will vary around that ’up’. But there is no foolproof solution here.}

\response We have expanded the explanation of weak identifiability: components with nearly equal rates can exchange weight with little change in the fitted density, while the rate of a component with negligible weight can be poorly constrained. These configurations can make covariance-based uncertainty propagation unreliable, motivating our use of bootstrap intervals.

We agree that bootstrap resampling reflects fluctuations around the observed dataset and does not automatically correct an upward fluctuation or background-modeling bias. In our study, however, resampling is used only to construct the confidence interval for each pseudo-dataset. Coverage is then evaluated across independently generated pseudo-datasets by checking whether those intervals contain the known data-generating density. This outer simulation directly assesses the performance of the bootstrap interval procedure, including any undercoverage arising from the concern raised by the referee. We have clarified this distinction in the manuscript.

\addedtext{Coverage is evaluated by calculating how often the true density falls within the estimated pointwise confidence intervals across independently generated pseudo-datasets. Constructing these intervals requires accounting for the joint dependence among fitted parameters and their nonlinear mapping to the density. When a local quadratic approximation to the likelihood is adequate, the estimated parameter covariance matrix can be used for approximate uncertainty propagation. In finite exponential mixtures, however, individual component parameters can be poorly constrained even when the combined density is comparatively well constrained. Components with nearly equal rates can exchange weight with little change in the density, while the rate of a component with negligible weight is weakly constrained by the data. These configurations can produce nearly flat likelihood directions, making the quadratic approximation unreliable. We therefore use the non-parametric bootstrap \cite{Efron_1986} to construct pointwise confidence intervals directly from an ensemble of fitted densities.

For each independently generated pseudo-dataset $D$, we draw $B$ bootstrap resamples of the same size by sampling its observations with replacement. Each candidate mixture order is refitted to every resample, starting from the parameter estimates obtained for that order by the random-restart procedure applied to $D$. The random-restart search is not repeated within each resample. The number of components $k$ is selected again using the same information criterion, so the ensemble includes variation associated with both parameter estimation and model selection. At each mass value $x$, the pointwise confidence interval is defined by the $100\alpha/2$ and $100(1-\alpha/2)$ percentiles of the fitted bootstrap densities. We use $\alpha=0.05$, corresponding to the 2.5th and 97.5th percentiles.

These intervals reflect variability estimated from each pseudo-dataset, but resampling does not automatically correct bias from background misspecification or model selection. Their nominal 95\% confidence level therefore does not guarantee 95\% coverage. To assess their actual coverage, we repeat the full interval-construction procedure across independently generated pseudo-datasets and record how often the resulting intervals contain the known true density.}

\refcomment{Metrics: you discuss coverage, spurious signal and Bias, all of which are
interesting, but you do not compare the power of the different approaches,
which is an unfortunate omission. AIC and BIC for example will surely differ
in power.}

\response We agree that power studies would strengthen the assessment, particularly by quantifying the potential loss of sensitivity associated with additional background flexibility. Completing and validating this comparison would require additional work beyond what we can reliably incorporate within the present revision period. We have therefore retained the scope of the current simulation studies and revised the discussion to state explicitly that the background-only results do not establish the preferred model complexity for a resonance search.

\refcomment{The Discussion is misleading. You had to abandon the basis functions drawn from EVT, so you should not now appeal to it to support your proposal.}
\response The positive-shape GPD remains part of our motivation: its exact representation as an exponential mixture with gamma-distributed rates provides the starting point for the proposed generalization. However, we agree that EVT does not establish that an arbitrary collider spectrum is an exponential mixture over a finite analysis range. We have revised the manuscript to distinguish this statistical motivation from the empirical assessment of the model. In particular, being far from the kinematic endpoint does not by itself justify an unbounded GPD approximation; its relevance here comes from its compatibility with the ATLAS and CMS spectra.

\refcomment{The rest of the discussion i applaud, bar the lack of discussion of sensitivity
- the price paid for multiple parameters. Without that there is really no
reason not to go to N → infinity. Without this you miss an important test.}
\response We agree, this should be in the discussion. The following text has been added:
\addedtext{
Additional background flexibility can reduce sensitivity to a resonant signal, although this tradeoff should be assessed alongside robustness to background misspecification. Exponential mixtures remain completely monotone as components are added, restricting their ability to reproduce localized resonant structure. These constraints do not eliminate partial signal absorption: when a localized deviation is compatible with an admissible variation in the background, increased uncertainty in the extracted signal yield can reflect the limited information available to distinguish these explanations. The background-only studies presented here favor the AIC-selected mixtures over the BIC-selected mixtures in terms of the bias, coverage, and spurious signal considered. These results support retaining additional background flexibility for density estimation, but do not establish the preferred complexity for a resonance search. Determining that choice requires assessing signal recovery and discovery power alongside control of false-positive rates.
}

\refcomment{Fig 2 and 4. Perhaps I have misunderstood something in your procedure?
As we take the limit of large x the exponential with the smallest weight beta
will increase in fraction. The w parameters are by construction non-negative.
So surely the log of the function has to be concave, which is not the case in
these figures. ? Is this not implicit in your monotone discussion? If negative
W’s are allowed there is of course more freedom..but also more fine tuning
to avoid negative PDFs.}
\response Figures 2 and 4 were plotted on a log-log scale. The exponential mixture is log-convex not log-concave. We have expanded discussion around log-convexity, as a result we have decided to use log plots instead of log-log plots.

\bigskip
\hrule
\bigskip

\section*{Response to Referee 2}

\refcomment{The lack of experiments with larger dataset sizes is unfortunate, as
most concern the rather small CMS diphoton dataset case. Did you
attempt simulation studies with the ATLAS one as well? A result
similar to the integrated luminosity scaling in ref [9] would serve as a
great motivation, as it stands this feels like a gap to me.}
\response We agree that this would be nice to have. Given the increased computational burden we have decided not to add this study. We have added the following text to the discussion:
\addedtext{Further simulation studies could more fully characterize the performance and applicability of the exponential mixture framework. Signal-injection and sensitivity studies could assess performance in the presence of a signal and evaluate how the choice of background model and its complexity affects discovery potential and exclusion limits. Studies at larger and smaller sample sizes would help establish whether the results obtained with 5,036 observations extend to other regimes. Although the fit to the observed dijet spectrum demonstrates applicability to a much larger dataset, it does not establish the corresponding repeated-sample performance. Extending the simulations to other background shapes would further clarify how broadly the observed performance holds.}

\refcomment{The paper describes background modeling and spurious signal studies,
but does not report on signal injection tests. Adding flexibility to
a background model can result in loss of power for a signal search
and corresponding experiments would be helpful. Both likelihood ratio
studies and linearity / extracted-vs-true signal tests would strengthen
this paper.}

\response We agree that signal-injection and power studies would strengthen the assessment, particularly by quantifying the potential loss of sensitivity associated with additional background flexibility. Completing and validating this comparison would require additional work beyond what we can reliably incorporate within the present revision period. We have therefore retained the scope of the current simulation studies and revised the discussion to state explicitly that the background-only results do not establish the preferred model complexity for a resonance search.

\refcomment{title page: arXiv identifier is a placeholder, strange line break before "USA"}
\response Fixed.

\refcomment{"extreme value theory" is sometimes capitalized and sometimes not}
\response Fixed, now using lower case in all instances.

\refcomment{sec 1: "Higgs mass occurs within the bulk of the background" is awkward phrasing, suggest "resonance" or similar}
\response Fixed, changed "Higgs mass" to "Higgs resonance".

\refcomment{sec 2: this is very statistics theory focused and the reader may think
that exponentials have not been used at all in practice, would be useful
to mention here that exponential sums are already common practice
(as the later examples show)}
\response Added text at end of introduction:
\addedtext{Exponential functions and sums of exponentials have been used previously for background modeling; here, we develop a statistical motivation for this model family through its connection to extreme value theory.}

\refcomment{sec 3.1: superfluous "be" in "expected to be require"}
\response Fixed.

\refcomment{sec 3.1 (and following): how do you convince yourself that 80 random
restarts and 100 retries is a sensible choice and that you get sufficiently
close to the desired (global) minimum? The discussion section hints at
robust optimization being an open question, do you have any practical
advice to incorporate here based on the experiments in this paper?}

\response The choices of restarts and retries are heuristic and better solutions exist. We have added some text to point to this limitation and to suggest improvements for future work.
\addedtext{The restart and fit-call budgets are specified separately for each example and should be regarded as heuristic implementation choices.}
\addedtext{Optimization remains a practical limitation of the present implementation. The likelihood is non-convex in the mixture parameters and can contain nearly flat directions when components have similar rates or negligible weights. Random restarts reduce dependence on initialization, but the finite search used here does not certify that the global maximum has been reached. The restart and retry budgets are heuristic, and their sufficiency across datasets and model complexities has not been established. Residual optimization error can also affect model selection when competing information-criterion values are close. An adaptive stopping criterion could help allocate computational effort according to the observed progress of the optimization, reducing reliance on manually chosen budgets while providing a reproducible rule for terminating the search. Beyond refining the current restart procedure, alternative strategies such as split-and-merge methods developed for Gaussian mixtures \cite{split-and-merge} warrant investigation for their potential to improve optimization reliability and efficiency in exponential mixtures.}

\refcomment{sec 3.1: "spanning from 2015 to 2018 ." extra space before period}
\response Fixed.

\refcomment{table 1: would suggest stating that the $\nu$ in $\chi_\nu^2$ is the number of degrees
of freedom (notation might not be obvious to all the HEP audience) and ideally also stating what that number is for completeness}
\response This has been added to the table captions.

\refcomment{fig 1: after table 1 I was trying to understand how the dijet function
enters here until I realized this is only showing exponential mixtures,
may help mentioning in the text introducing the figure explicitly that
this is only showing those}
\response The text has been updated:
\addedtext{Figure~\ref{fig:dijet_AIC_BIC} compares the AIC and BIC scores for exponential mixtures of varying complexity to determine the best fit for the dijet dataset. In this case AIC and BIC agree, the best model has four components. The goodness-of-fit metrics are shown in Table~\ref{tab:dijet_AIC_BIC} where we compare between the dijet function and the exponential mixture. The exponential mixture achieves a slightly smaller AIC and $\chi^2$ but uses more parameters.
Figure~\ref{fig:dijet_fit} shows the fit of the dijet function and the 4-component exponential mixture to the dijet dataset.
The residuals show that the two functions provide similar estimates throughout the mass range.}

\refcomment{fig 1 and fig 3: would be nice to replace with a vector graphic version
if available, though perfectly readable at current resolution}
\response These figures have been replaced with vector graphic versions.

\refcomment{fig 5: legend placement overlaps a bit with the figure content, looks
readable enough to me but perhaps can be tuned slightly}
\response The legend text has been made slightly smaller.

\refcomment{sec 4: "the functions in eq. (3.4)" should be 3.4 to 3.7 (also again later)}
\response Fixed, these functions are now labeled 3.4a,3.4b, ...

\refcomment{sec 4: I think a brief explanation of label switching is useful here for a
more general audience}
\response We have replaced the brief reference to label switching with an explicit description of weak identifiability. Permuting component labels produces equivalent parameterizations of the same density. Separately, nearly coincident rates or negligible component weights can leave individual parameters poorly constrained, producing nearly flat likelihood directions. The revised text focuses on this latter issue because it is directly relevant to uncertainty estimation.

\addedtext{In finite exponential mixtures, individual component parameters can be poorly constrained even when the combined density is comparatively well constrained. For example, components with nearly equal rates can exchange weight with little change in the density, while the rate of a component with negligible weight is weakly constrained by the data. These configurations can produce nearly flat likelihood directions, making a local quadratic approximation to the likelihood unreliable.}

\refcomment{sec 4: how trustworthy is the Hessian-based estimate of $\sigma_s$ in the spurious signal section, given that the previous coverage section discusses the Hessians are unstable?}
\response This is a good question, for most of the fits the Hessian is forced positive definite even after model selection. Instead of quantifying the spurious signal in terms of $N_s/\sigma_{N_s}$ we have decided to use the signed profile likelihood ratio which is equivalent to this quantity under the Wald approximation. See the added text:
\addedtext{We quantify the spurious signal using the signed profile likelihood statistic
\begin{equation}
r_0 = \operatorname{sign}(\hat{\mu})\sqrt{-2\ln\lambda(0)},
\end{equation}
where $\lambda(\mu)$ is the profile likelihood ratio defined in Eq.~(7) of Ref.~\cite{Cowan_2011}. Its sign distinguishes apparent excesses from deficits. Under the Wald approximation,
$-2\ln\lambda(0)\approx\hat{\mu}^2/\sigma_\mu^2$, so $r_0$ reduces to the commonly used ratio $\hat{\mu}/\sigma_\mu$ (see Eqs.~(17)--(18) of Ref.~\cite{Cowan_2011}). The uncertainty $\sigma_\mu$ is typically estimated from the inverse Hessian of the negative log likelihood, whose inversion can be unstable when the likelihood is flat. Using $r_0$ avoids this explicit covariance-matrix inversion and does not require a local quadratic approximation to the likelihood, although its interpretation as a Gaussian significance still relies on asymptotic approximations.
}

\refcomment{sec 4: typo "regio;,"}
\response Fixed.

\refcomment{fig 7: caption is about S(x) but figure is $N_s/\sigma_{N_s}$}
\response Fixed.

\refcomment{sec 5: EVT abbreviation is re-defined here}
\response Fixed.

\refcomment{sec 5: stray "due;"}
\response Fixed.

\refcomment{ref [25] uses a different title formatting compared to the rest of the references}
\response Ref. 25 was a reference to the HEPData entry; this has been removed since a link to the HEPData entry is provided in the data availability section.

\refcomment{ref [25] uses "Collaboration", while the rest uses "collaboration"}
\response Fixed.

\bigskip
\hrule
\bigskip

\noindent We hope the revised manuscript is now suitable for publication in JHEP.

\bigskip
\noindent Sincerely,\\
\textit{Austin Townsend, Marc Osherson, Mike Hildreth, Stefano Castruccio}

\end{document}